\chapter{Literature Review}
\label{ch:literature_review}

The study of transverse momentum ($p_T$) spectra in high-energy proton-proton and heavy-ion collisions has a rich history, with phenomenological models continually evolving to capture both the soft (thermal) and hard (perturbative QCD) regimes of particle production. This chapter reviews the foundational models and recent advancements that form the baseline for our validation framework.

\section{Thermal and Hydrodynamical Foundations}
Early approaches to modeling particle spectra relied heavily on the Boltzmann-Gibbs statistics, assuming local thermal equilibrium of the produced fireball. However, the observation of power-law tails at high $p_T$ necessitated refinements. The Blast-Wave model \citep{Schnedermann1993} introduced collective transverse flow, successfully describing the soft bulk of the spectra ($p_T \lesssim 2.5$ GeV/c) by assuming a locally thermalized source expanding with a collective transverse velocity profile. 

The standard Boltzmann-Gibbs Blast-Wave (BGBW) model evaluates the invariant yield as an integral over the freeze-out surface. While highly successful at low $p_T$, the BGBW model systematically underpredicts the yield at higher transverse momenta, where hard parton scattering dominates.

\section{The Tsallis Distribution and Non-Extensive Thermodynamics}
To simultaneously describe both the low-$p_T$ thermal bulk and the high-$p_T$ power-law tail, the non-extensive Tsallis statistics were introduced to high-energy physics. A seminal advancement was the formulation of a thermodynamically consistent Tsallis distribution by \citet{Cleymans2012}, which strictly satisfies thermodynamic relations for particle number, energy density, and pressure. The Tsallis distribution introduces a non-extensivity parameter $q$; as $q \to 1$, the distribution recovers the standard Boltzmann-Gibbs form.

\section{Recent Advancements and Multiplicity Dependence}
Primary measurements of charged-particle $p_T$ spectra at the LHC have been extensively provided by the ALICE Collaboration \citep{ALICE2018}. These foundational datasets spanning multiple collision systems ($pp$, $p$-Pb, and Pb-Pb) and collision energies (up to $\sqrt{s} = 13$ TeV) serve as the empirical bedrock for phenomenological modeling.

Recent studies have focused heavily on the multiplicity dependence of freeze-out parameters in these proton-proton collisions. \citet{Khuntia2019} analyzed $p_T$ spectra at $\sqrt{s} = 7$ TeV using both the BGBW and the thermodynamically consistent Tsallis distribution. They observed a mass ordering in the kinetic freeze-out temperature $T_{\text{kin}}$ and the non-extensivity parameter, supporting a differential freeze-out scenario. Crucially, their BGBW analysis showed that the average radial flow velocity $\langle \beta \rangle$ remains largely independent of the event multiplicity, providing a key benchmark for multi-component models.

\citet{Rath2020} extended this work to 5.02 and 13 TeV collision energies. Their findings corroborated that $\langle \beta \rangle$ is almost independent of collision energy and multiplicity. However, they found that $T_{\text{kin}}$ significantly depends on the multiplicity classes. They concluded that while the Tsallis model better describes the full $p_T$ range, the BGBW model remains highly effective for the bulk part ($p_T \lesssim 2.5$ GeV/c).

More recently, \citet{Biro2025} proposed a two-component soft/hard baseline model for ALICE $pp$ data from 2.76 to 13 TeV. Their analysis revealed that a simple Boltzmann fit effectively describes the soft component across all energies, while the residual hard spectra exhibit no evolution in shape or peak position with event multiplicity. Furthermore, the mean $p_T$ for both the soft and hard components remains nearly constant across multiplicity classes, challenging more complex hydrodynamical interpretations that necessitate highly parameterized models.

These literature benchmarks underscore the necessity of rigorous statistical validation when proposing novel multi-component models, such as the 3-component formulation (incorporating a moving J\"uttner distribution) evaluated in this thesis.
